\section*{Abstract}
Multi-agent systems (MAS) exhibit complex emergent behavior that has pushed the frontier of large language model
(LLM) capabilities. In this work, we examine how these agent swarms respond to positive (regulators) and negative
(disregulators, poisoned) peer pressure in tasks that elicit reward hacking. We integrate and adapt
\textit{CheatBench} tasks into the \textit{ORBIT} framework to run experiments that measure how reward-hacking and
collaboration in swarms are affected by i) the proportion of regulator and disregulator agents, ii) the onset and
active window of peer pressure, and iii) the strength and ambiguity of instruction. Our investigation is ongoing.
Our preliminary results show that regulators can reduce cheating in software engineering (SWE) and bioresearch
evaluations, suggesting that positive peer pressure is an plausible mitigation technique for emergent cheating and
collusion, but that effectiveness is often limited by limited proactive communication of regulators. A regulator
that warned the team early cut cheating on SWE. Disregulators raised attempts at cheating but led no team to use the
planted work.
